\documentclass[10pt,a4paper]{amsart}
\usepackage[margin=19mm]{geometry}
\usepackage{amsmath,amssymb,mathtools}
\usepackage[scr=rsfs,cal=euler]{mathalfa}
\usepackage{xfrac}
\usepackage[unicode,hidelinks,bookmarks=false]{hyperref}
\newcommand{\ie}{i.e.\ }
\newcommand{\cf}{cf.\ }
\newcommand{\resp}{respectively\ }
\newcommand{\Subsection}{\subsection}
\providecommand{\nobreakdash}{\nobreak\hbox{-}}
\setlength{\emergencystretch}{2em}
\multlinegap0pt
\usepackage{bm}
\usepackage{bbm}
\usepackage{dsfont}
\usepackage{xspace}
\usepackage{cancel}
\usepackage{mathtools}
\usepackage{amsthm}
\makeatletter
\@ifundefined{stat}{\newtheorem{stat}{Stat}}{}
\@ifundefined{lemma}{\newtheorem{lemma}[stat]{Lemma}}{}
\makeatother
\newcommand{\E}{\mathrm{E}}
\newcommand{\N}{\mathbb{N}}
\newcommand{\R}{\mathbb{R}}
\newcommand{\cC}{\mathcal{C}}
\newcommand{\e}{{\rm e}}
\renewcommand{\d}{{\rm d}}
\renewcommand{\ge}{\geqslant}
\renewcommand{\le}{\leqslant}
\title[On the passage times of self-similar Gaussian processes on c]{On the passage times of self-similar Gaussian processes on curved boundaries}
\author{Davar Khoshnevisan, Cheuk Yin Lee}
\date{Source: arXiv:2506.15949v2 (CC BY 4.0)}
\begin{document}
\maketitle

\noindent\emph{Source and licence.} On the passage times of self-similar Gaussian processes on curved boundaries. Version arXiv:2506.15949v2 (CC BY 4.0), distributed under the Creative Commons Attribution 4.0 International licence, as recorded in the arXiv licence metadata for this exact version and shown on the abstract page https://arxiv.org/abs/2506.15949v2. Full rights notice: licenses/arxiv-2506.15949-CC-BY-4.0.txt.\par\smallskip

\noindent\textit{Editorial scope.} Counted unit: the Lemma whose source label is \texttt{lem:S:cont} in the article (source0.tex lines 1295--1297, source label \texttt{lem:S:cont}), delivered together with the article's own complete original proof of it (source0.tex lines 1299--1401, one \texttt{proof} environment of 103 lines). No mathematics is written, repaired or re-derived: statement and proof are the article's own text line for line. Required context retained uncounted: S (lines 1288--1290). Every definition, display and label of those lines is retained unchanged. Omitted from this package and not redistributed: 1--1287, 1291--1294, 1298--1298, 1402--1691. Nothing inside a retained excerpt is omitted or altered: every retained body is delivered exactly as the article's lines are written, so no omission receipt of any kind accompanies this package. Citations: the retained excerpts carry no citation command at all, so no bibliography entry of the article is transcribed and no reference is redistributed. Reference adaptation: none was needed and none was made; every reference inside the delivered unit resolves inside it, so no unresolved reference or citation is produced. Printed slips kept literal: no printed word, symbol, hypothesis or number of the counted statement or of its proof was altered. Layout and class: the article's own class is \texttt{amsart} with packages including graphicx, xcolor, cite, mathtools, bbold, bbding, amssymb, amsmath, amsthm, mathrsfs, paralist, bm, esint, setspace; the wrapper is the lane's pinned \texttt{amsart} editorial wrapper at 10pt on A4, which restates the theorem-like environments the retained text needs and adds the article's own macro definitions that the retained text needs (\texttt{\textbackslash E}, \texttt{\textbackslash N}, \texttt{\textbackslash R}, \texttt{\textbackslash cC}, \texttt{\textbackslash d}, \texttt{\textbackslash e}, \texttt{\textbackslash ge}, \texttt{\textbackslash le}), with extra wrapper packages (bm, bbm, dsfont, xspace, cancel, mathtools). The delivered numbering is the unit's own. Assets: no retained span contains a figure environment, a graphics-inclusion command, a PDF-page inclusion, a table or a TikZ picture; no image file is copied into this package, the package declares no asset dependency and no image file is redistributed with it. Licence: version arXiv:2506.15949v2 of this article is distributed under the Creative Commons Attribution 4.0 International licence, as recorded in the arXiv licence metadata for this exact version and shown on the abstract page https://arxiv.org/abs/2506.15949v2. A full rights notice is delivered at licenses/arxiv-2506.15949-CC-BY-4.0.txt. Characters: every retained excerpt is pure ASCII, so the delivered engine, XeLaTeX with its default Latin Modern text font, reports no missing character.
\begin{align}\label{S}
	S(t) = \sqrt{\frac{c(d\,,\beta)}{2(2\pi)^{d}}} \int_{\R^d} \frac{1-\e^{-t\|\xi\|^\gamma}}{\|\xi\|^{(d+\gamma-\beta)/2}} W_1(\d\xi) \qquad \forall t \ge 0,
\end{align}

\begin{lemma}\label{lem:S:cont}
	The process $S$ is a.s.~continuous on $\R_+$ and $\cC^\infty$ on $(0\,,\infty)$.
\end{lemma}

\begin{proof}
	For every $t,h>0$,
	\begin{align*}
		\E\big( |S(t+h)-S(t)|^2\big) 
			&\propto \int_{\R^d} \e^{-2t\|\xi\|^\gamma}
			\left( 1 - \e^{-h\|\xi\|^\gamma} \right)^2 \|\xi\|^{\beta-\gamma-d}\, \d \xi\\
		&\lesssim h^2 \int_{\R^d}  \e^{-2t\|\xi\|^\gamma} \|\xi\|^{\beta+\gamma-d}\, \d \xi
			\propto \frac{h^2}{t^{1+(\beta/\gamma)}},
	\end{align*}
	where the implied constants depend only on $(d\,,\beta\,,\gamma)$.
	In particular, 
	\begin{equation}\label{S:tmpsup}
		\sup_{t\ge h}\E\big( |S(t+h)-S(t)|^2\big) \lesssim h^{1-(\beta/\gamma)}\quad
		\forall h>0,
	\end{equation}
	where the implied constant depends only on $(d\,,\beta\,,\gamma)$.
	Also, by \eqref{S} and scaling, $\E( |S(t)|^2)\propto t^{1-(\beta/\gamma)}$
	uniformly for every $t>0$. Therefore,
	\begin{align*}
		\E\big( |S(t+h)-S(t)|^2\big) &\le 2\E\big(|S(t+h)|^2\big)+2\E\big(|S(t)|^2\big)\\
		&\lesssim (t+h)^{1-(\beta/\gamma)}+ t^{1-(\beta/\gamma)}
			\qquad\forall t,h\ge0.
	\end{align*}
	This and \eqref{S:tmpsup} together yield
	\[
		\E\big( |S(t+h)-S(t)|^2\big) \lesssim h^{1-(\beta/\gamma)}
		\quad\forall t,h>0,
	\]
	where the implied constant depends only on $(d\,,\beta\,,\gamma)$. Thanks to the Kolmogorov
	continuity theorem and the fact that moments of Gaussian random variables
	are determined by their variance, it follows that that $S\in\cC^a_{\textit{loc}}(\R_+)$
	a.s.\ for every $a\in(0\,, \frac12(1-\beta\gamma^{-1}))$. In particular, $S$
	is continuous on $\R_+$. Next define
	\[
		S^{(n)}(t)
		= \sqrt{\frac{c(d\,,\beta)}{2(\pi)^d}}(-1)^{n+1} t^n \int_{\R^d} 
		\|\xi\|^{\frac12(\beta-d)+(n-\frac12)\gamma}\e^{-t\|\xi\|^\gamma}
		\,W_1(\d\xi) \quad\forall t > 0,\ n\in\N,
	\]
	where $W_1$ is the same 
	white noise that was used in \eqref{S} to define $S$. The process $S^{(n)}$ is well defined,
	for example because
	\begin{align*}
		\E\left(|S^{(n)}(t)|^2\right) \propto \int_{\R^d} 
		\|\xi\|^{\beta-d+(2n-1)\gamma} \e^{-2t\|\xi\|^\gamma}\,\d\xi
		<\infty\quad\forall t > 0,\ n\in\N.
	\end{align*}
	Define
	\[
		k(t\,,\xi) = 
		\sqrt{\frac{c(d\,,\beta)}{2(2\pi)^d}} \left(\frac{ 1 - \e^{-t\|\xi\|^\gamma}}{%
		\|\xi\|^{(d+\gamma-\beta)/2}}\right)\qquad\forall t>0,\
		\xi\in\R,
	\]
	and observe that 
	\[
		S(t)=\int_{\R^d} k(t\,,\xi)\,W_1(\d\xi),\quad
		S^{(n)}(t) = \int_{\R^d} \partial^n_t k(t\,,\xi)\, W_1(\d\xi).
	\]
	Thanks to the preceding $L^2$-computations, we may apply a stochastic Fubini
	argument twice in order to see that, for all non random $f\in\cC^\infty_c(0\,,\infty)$,
	\begin{align*}
		&\int_0^\infty S^{(n)}(t)f(t)\,\d t = \int_{\R^d} W_1(\d\xi)
			\int_0^\infty\d t\ f(t)\partial^n_t k(t\,,\xi)\\
		&= (-1)^n\int_{\R^d} W_1(\d\xi)
			\int_0^\infty\d t\ \frac{\d^n f(t)}{\d t^n} k(t\,,\xi)
			= (-1)^n\int_0^\infty S(t)\frac{\d^n f(t)}{\d t^n}\,\d t
			\quad\text{a.s.}
	\end{align*}
	It follows that $S^{(n)}$ is the $n$th weak derivative of $S$, and hence it remains to 
	prove that $S^{(n)}$ is continuous [up to a modification]. With that end in mind note that
	\[
		\E\left(|S^{(n)}(t) - S^{(n)}(s)|^2\right)\lesssim Q_1+Q_2,
	\]
	where the implied constant is universal, and
	\begin{align*}
		Q_1&= (t^n-s^n)^2 \int_{\R^d} \|\xi\|^{\beta-d+(2n-1)\gamma} \e^{-2t\|\xi\|^\gamma}\d\xi,\\
		Q_2&= s^{2n}\int_{\R^d} \|\xi\|^{\beta-d+(2n-1)\gamma}
			\left( \e^{-s\|\xi\|^\gamma} - \e^{-t\|\xi\|^\gamma}\right)^2\,\d\xi.
	\end{align*}
	Choose and fix two arbitrary numbers $\varepsilon$ and $M$ that satisfy
	$0<\varepsilon<1<M$. Since
	$t^n-s^n=n^{-1}\int_s^t x^{n-1}\,\d x$,
	\[
		Q_1 \le n^{-2}M^{2n-2}(t-s)^2\int_{\R^d} 
		\|\xi\|^{\beta-d+(2n-1)\gamma} \e^{-2\varepsilon\|\xi\|^\gamma}\d\xi
		\lesssim (t-s)^2,
	\]
	uniformly for all $s,t\in[\varepsilon\,,M]$. Also,
	\begin{align*}
		Q_2&\le M^{2n}\int_{\R^d} \|\xi\|^{\beta-d+(2n-1)\gamma}
			\left( 1-\e^{-|t-s| \|\xi\|^\gamma}\right)^2 
			\e^{-2\varepsilon\|\xi\|^\gamma}\d\xi\\
		&\le M^{2n}(t-s)^2\int_{\R^d} \|\xi\|^{\beta-d+(2n+1)\gamma}
			 \e^{-2\varepsilon\|\xi\|^\gamma}\d\xi \propto (t-s)^2,
	\end{align*}	
	uniformly for all $s,t\in[\varepsilon\,,M]$.
	The implied constants above depend only on 
	$(d\,,\beta\,,\gamma\,,n\,,\varepsilon\,,M)$.
	Thus, we can deduce from Kolmogorov's continuity theorem 
	that $S^{(n)}$ is continuous on $(0\,,\infty)$. 
	 This concludes the proof of Lemma \ref{lem:S:cont}.
\end{proof}

\end{document}
